%%%%%%%%%%%%%%%%%%%%%%%%%%%%%%%%%%%%%%%%%
% Masters/Doctoral Thesis 
% LaTeX Template
% Version 1.43 (17/5/14)
%
% This template has been downloaded from:
% http://www.LaTeXTemplates.com
%
% Original authors:
% Steven Gunn 
% http://users.ecs.soton.ac.uk/srg/softwaretools/document/templates/
% and
% Sunil Patel
% http://www.sunilpatel.co.uk/thesis-template/
%
% License:
% CC BY-NC-SA 3.0 (http://creativecommons.org/licenses/by-nc-sa/3.0/)
%
% Note:
% Make sure to edit document variables in the Thesis.cls file
%
%
% Modified and Adapted by Michael Lees 2020
%%%%%%%%%%%%%%%%%%%%%%%%%%%%%%%%%%%%%%%%%

%----------------------------------------------------------------------------------------
%	PACKAGES AND OTHER DOCUMENT CONFIGURATIONS
%----------------------------------------------------------------------------------------

\documentclass[11pt, oneside]{Thesis} % The default font size and one-sided printing (no margin offsets)

\graphicspath{{Pictures/}} % Specifies the directory where pictures are stored

\usepackage[square, numbers, comma, sort&compress]{natbib} % Use the natbib reference package - read up on this to edit the reference style; if you want text (e.g. Smith et al., 2012) for the in-text references (instead of numbers), remove 'numbers' 
\usepackage[ruled,vlined]{algorithm2e}
\usepackage{amsmath}
\usepackage{booktabs}
\DeclareMathOperator*{\argmin}{arg\,min}
\DeclareMathOperator*{\argmax}{arg\,max}
\hypersetup{urlcolor=blue, colorlinks=true} % Colors hyperlinks in blue - change to black if annoying
\title{\ttitle} % Defines the thesis title - don't touch this

\begin{document}

\frontmatter % Use roman page numbering style (i, ii, iii, iv...) for the pre-content pages
\setstretch{1.3} % Line spacing of 1.3

% Define the page headers using the FancyHdr package and set up for one-sided printing
\fancyhead{} % Clears all page headers and footers
\rhead{\thepage} % Sets the right side header to show the page number
\lhead{} % Clears the left side page header

\pagestyle{fancy} % Finally, use the "fancy" page style to implement the FancyHdr headers

\newcommand{\HRule}{\rule{\linewidth}{0.5mm}} % New command to make the lines in the title page

% PDF meta-data
\hypersetup{pdftitle={\ttitle}}
\hypersetup{pdfsubject=\subjectname}
\hypersetup{pdfauthor=\authornames}
\hypersetup{pdfkeywords=\keywordnames}

%----------------------------------------------------------------------------------------
%	TITLE PAGE
%----------------------------------------------------------------------------------------

\begin{titlepage}
\begin{center}

\textsc{\LARGE \univname}\\[1.5cm] % University name
\textsc{\Large Masters Thesis}\\[0.5cm] % Thesis type

\HRule \\[0.4cm] % Horizontal line
{\huge \bfseries \ttitle}\\[0.4cm] % Thesis title
\HRule \\[1.5cm] % Horizontal line
 
\begin{minipage}{0.4\textwidth}
\begin{flushleft} \large
\emph{Author:}\\
{\authornames} % Author name
\end{flushleft}
\end{minipage}
\begin{minipage}{0.4\textwidth}
\begin{flushright} \large
\emph{Supervisors:} \\
{\supname} % Supervisor name
\end{flushright}
\end{minipage}\\[3cm]

\large \textit{A thesis submitted in partial fulfilment of the requirements\\ for the degree of \degreename}\\[0.3cm] % University requirement text
\textit{in the}\\[0.4cm]
\groupname\\\deptname\\[2cm] % Research group name and department name
 
{\large \today}\\[2cm] % Date
\includegraphics[width=0.6\textwidth]{Figures/clslogo.png} % Include Computational Science Logo
 
\vfill
\end{center}

\end{titlepage}

%----------------------------------------------------------------------------------------
%	DECLARATION PAGE
%	Your institution may give you a different text to place here
%----------------------------------------------------------------------------------------

\Declaration{

\addtocontents{toc}{\vspace{1em}} % Add a gap in the Contents, for aesthetics

I, \authornames, declare that this thesis, entitled `\ttitle' and the work presented in it are my own. I confirm that:

\begin{itemize} 
\item[\tiny{$\blacksquare$}] This work was done wholly or mainly while in candidature for a research degree at the University of Amsterdam.
\item[\tiny{$\blacksquare$}] Where any part of this thesis has previously been submitted for a degree or any other qualification at this University or any other institution, this has been clearly stated.
\item[\tiny{$\blacksquare$}] Where I have consulted the published work of others, this is always clearly attributed.
\item[\tiny{$\blacksquare$}] Where I have quoted from the work of others, the source is always given. With the exception of such quotations, this thesis is entirely my own work.
\item[\tiny{$\blacksquare$}] I have acknowledged all main sources of help.
\item[\tiny{$\blacksquare$}] Where the thesis is based on work done by myself jointly with others, I have made clear exactly what was done by others and what I have contributed myself.\\
\end{itemize}


Signed: 

%\\\\\includegraphics[width=3cm]{Pictures/yoursignature.pdf} \\
\rule[1em]{25em}{0.5pt} % This prints a line for the signature
 
Date: DATE
\rule[1em]{25em}{0.5pt} % This prints a line to write the date
}

\clearpage % Start a new page

%----------------------------------------------------------------------------------------
%	QUOTATION PAGE
%----------------------------------------------------------------------------------------

\pagestyle{empty} % No headers or footers for the following pages

\null\vfill % Add some space to move the quote down the page a bit

\textit{``What I cannot create, I do not understand."}

\begin{flushright}
Richard P. Feynman
\end{flushright}

\vfill\vfill\vfill\vfill\vfill\vfill\null % Add some space at the bottom to position the quote just right

\clearpage % Start a new page

%----------------------------------------------------------------------------------------
%	ABSTRACT PAGE
%----------------------------------------------------------------------------------------

\addtotoc{Abstract} % Add the "Abstract" page entry to the Contents

\abstract{\addtocontents{toc}{\vspace{1em}} % Add a gap in the Contents, for aesthetics

Synthetic data, as the name suggests, is data which is artificially created rather than being generated by actual events. It’s often generated by algorithms or by AI methods and can be used as test data for new products or for AI model training. 
The main benefit of synthetic data is that it can be generated to meet specific conditions that are not available in (real) data. In this thesis the aim is to improve imbalanced learning algorithms by synthesizing new data to aid model learning. Current methods are difficult to train and rely solely on machine learning techniques which makes it less intuitive.
This novel approach was used to aid in classification problems and compared to other oversampling techniques
}

\clearpage % Start a new page

%----------------------------------------------------------------------------------------
%	ACKNOWLEDGEMENTS
%----------------------------------------------------------------------------------------

\setstretch{1.3} % Reset the line-spacing to 1.3 for body text (if it has changed)

\acknowledgements{\addtocontents{toc}{\vspace{1em}} % Add a gap in the Contents, for aesthetics

Thank the people that have helped, supervisors family etc.

}
\clearpage % Start a new page

%----------------------------------------------------------------------------------------
%	LIST OF CONTENTS/FIGURES/TABLES PAGES
%----------------------------------------------------------------------------------------

\pagestyle{fancy} % The page style headers have been "empty" all this time, now use the "fancy" headers as defined before to bring them back

\lhead{\emph{Contents}} % Set the left side page header to "Contents"
\tableofcontents % Write out the Table of Contents

\lhead{\emph{List of Figures}} % Set the left side page header to "List of Figures"
\listoffigures % Write out the List of Figures

\lhead{\emph{List of Tables}} % Set the left side page header to "List of Tables"
\listoftables % Write out the List of Tables

\lhead{\emph{List of Algorithms}} % Set the left side page header to "List of Algorithms"
\addtotoc{List of Algorithms}
\listofalgorithms % Write out the List of Tables

%----------------------------------------------------------------------------------------
%	ABBREVIATIONS
%----------------------------------------------------------------------------------------

\clearpage % Start a new page

\setstretch{1.5} % Set the line spacing to 1.5, this makes the following tables easier to read

\lhead{\emph{Abbreviations}} % Set the left side page header to "Abbreviations"
\listofsymbols{ll} % Include a list of Abbreviations (a table of two columns)
{
\textbf{CSL} & \textbf{C}omputational \textbf{S}ceince \textbf{L}ab\\
\textbf{UvA} & \textbf{U}niversitiet \textbf{v}an \textbf{A}msterdam\\
%\textbf{Acronym} & \textbf{W}hat (it) \textbf{S}tands \textbf{F}or \\
}

%----------------------------------------------------------------------------------------
%	PHYSICAL CONSTANTS/OTHER DEFINITIONS
%----------------------------------------------------------------------------------------

\clearpage % Start a new page

%\lhead{\emph{Physical Constants}} % Set the left side page header to "Physical Constants"

%\listofconstants{lrcl} % Include a list of Physical Constants (a four column table)
%{
%Speed of Light & $c$ & $=$ & $2.997\ 924\ 58\times10^{8}\ \mbox{ms}^{-\mbox{s}}$ (exact)\\
%% Constant Name & Symbol & = & Constant Value (with units) \\
%}

%----------------------------------------------------------------------------------------
%	SYMBOLS
%----------------------------------------------------------------------------------------

%\clearpage % Start a new page

%\lhead{\emph{Symbols}} % Set the left side page header to "Symbols"

%\listofnomenclature{lll} % Include a list of Symbols (a three column table)
%{
%$a$ & distance & m \\
%$P$ & power & W (Js$^{-1}$) \\
%% Symbol & Name & Unit \\

%& & \\ % Gap to separate the Roman symbols from the Greek
%
%$\omega$ & angular frequency & rads$^{-1}$ \\
%% Symbol & Name & Unit \\
%}

%----------------------------------------------------------------------------------------
%	DEDICATION
%----------------------------------------------------------------------------------------

\setstretch{1.3} % Return the line spacing back to 1.3

\pagestyle{empty} % Page style needs to be empty for this page

\addtocontents{toc}{\vspace{2em}} % Add a gap in the Contents, for aesthetics

Deze is voor alle drerries

%----------------------------------------------------------------------------------------
%	THESIS CONTENT - CHAPTERS
%----------------------------------------------------------------------------------------

\mainmatter % Begin numeric (1,2,3...) page numbering

\pagestyle{fancy} % Return the page headers back to the "fancy" style

% Include the chapters of the thesis as separate  tex files from the Chapters folder
% Change the file names if you prefer

%This structure provides a bare bones essentials of the thesis and some indicative length
%This structure may not fit your thesis perfectly, but be sure to include these components somehow.
%It is possible to split the chapters up (E.g., 2 methods chapter, Experiment and Results as two, etc.)

% The typical length may be between 43 - 64 pages. Do not worry if you go slightly larger or smaller than this. 
% But a thesis of 20 pages, or 100 pages may suggest you've been to brief or too verbose.
%NOTE - there is not strict limit/minimum.
\chapter{Introduction}
\lhead{\emph{Introduction}}

Imbalanced learning is a term used to describe certain classification problems. In a classification problem the goal is to train a classifier to distinguish samples from different classes. In practice this could mean a classifier which determines if a patient has a disease or not, or more finance related, whether a transaction is fraudulent or not.

The imbalance implies that there is a mismatch between the sizes of the classes, something that is very common in real data sets. In the example of fraud detection: fraud transactions only make up a minute amount of all transactions. This leads to a problem when training a classifier. A classifier assumes that the classes are (roughly) equal in size and will create a model which distinguishes between the two classes. When one class significantly outweighs the other and the decision boundary becomes unclear. As a result, the classifier will favour the majority class and 

However, most classifiers will still appear to score well even 

\section{Related work}

oversamplers %Set out your thesis and state your research question (5-8 pages)

\chapter{Theory}
\lhead{Theory}
The literature discussed in the previous chapter provides a base for Synthsonic. While SMOTE is easy to use and GANs are very successful at generating images and sound, there is still room for a probabilistic model for synthesizing tabular data due to the limitations of regular SMOTE-based oversamplers and the difficulties of GANs.

\section{Synthsonic}
This section is dedicated to the theory behind Synthsonic. Contrary to other techniques such as SMOTE or other generative algorithms, Synthsonic uses reversible mathematical transformations, making use of as much information in the dataset as possible. This has the additional benefit that Synthsonic is easier to interpret than other GANs, where the synthesis of new samples happens through convolutional neural networks. Additionally, 


Synthsonic consists out of multiple reversible transformations of features. It can handle both numerical and categorical features, and will split a dataset $X$ into a numerical and categorical set, so that $X = (X_{num}, X_{cat})$. Modelling of the dataset happens in five steps:

\begin{enumerate}
    \item Numerical features are transformed into uniformly distributed ones, using reversible steps, from which the first-order correlations have been removed.
    \item Categorical and discretized numerical features are jointly modelled with a Bayesian network.
    \item Specific features can be selected for additional modelling.
    \item A discriminative learner is trained to model the residual discrepancies observed between the input data and the Bayesian network.
    \item The discriminative learner is calibrated to reweigh samples from the Bayesian network.
\end{enumerate}

Each of these steps will be discussed over the following sections.

\section{}

\chapter{Methods}
\lhead{\emph{Methods}}

This section describes the methods used in this thesis. Several terms and metrics will be introduced to the reader which are important to understand the experiments and results. Firstly, an imbalanced dataset of size $N$ is divided into a minority $N_{min}$ and majority $N_{maj}$ part, based on the number of samples per class, where $N = N_{min} + N_{maj}$. In this thesis, the minority class may also be named as the positive class. The experiments investigate in what ways oversampling can improve a classifier score.  

\section{datasets}
In order to examine the performance of the oversamplers, 27 datasets are used from the imbalanced-learn package~\cite{Lemaitre2017Imbalanced-learn:Learning}. Almost all of these datasets are publicly available in the UCI Machine Learning Repository~\cite{Dua2017UCIRepository}, which offers a wide variety of datasets for various machine learning purposes. The Imbalanced-learn package provides a selected set of imbalanced datasets to systematically benchmark the performance of the oversamplers. The datasets are also binarized, resulting in two classes to predict. This means that datasets which are multiclass, such as 'wine\_quality', are combined in such a way that only two remain with an imbalanced amount of samples.

\begin{table}
    \centering
        \begin{tabular}{llrrrrr}
        \toprule
                {}& dataset &    size &  features &  numerical &  categorical & imbalance \\
                & & & & features & features & ratio \\
        \midrule
        1 &           ecoli &     336 &         7 &                   5 &                     2 &           8.6:1 \\
        2 & optical\_digits &    5620 &        64 &                   0 &                    64 &           9.1:1 \\
        3 &       satimage &    6435 &        36 &                  36 &                     0 &           9.3:1 \\
        4 &     pen\_digits &   10992 &        16 &                  16 &                     0 &           9.4:1 \\
        5 &        abalone &    4177 &        10 &                   7 &                     3 &           9.7:1 \\
        6 & sick\_euthyroid &    3163 &        42 &                   6 &                    36 &           9.8:1 \\
        7 &   spectrometer &     531 &        93 &                  93 &                     0 &          10.8:1 \\
        8 &    car\_eval\_34 &    1728 &        21 &                   0 &                    21 &          11.9:1 \\
        9 &         isolet &    7797 &       617 &                 610 &                     7 &          12.0:1 \\
        10 &       us\_crime &    1994 &       100 &                  99 &                     1 &          12.3:1 \\
        11 &     yeast\_ml8 &    2417 &       103 &                 103 &                     0 &          12.6:1 \\
        12 &          scene &    2407 &       294 &                 294 &                     0 &          12.6:1 \\
        13 &    libras\_move &     360 &        90 &                  90 &                     0 &          14.0:1 \\
        14 &   thyroid\_sick &    3772 &        52 &                   6 &                    46 &          15.3:1 \\
        15 &      coil\_2000 &    9822 &        85 &                   1 &                    84 &          15.8:1 \\
        16 &     arrhythmia &     452 &       278 &                 137 &                   141 &          17.1:1 \\
        17 & solar\_flare\_m0 &    1389 &        32 &                   0 &                    32 &          19.4:1 \\
        18 &            oil &     937 &        49 &                  39 &                    10 &          21.9:1 \\
        19 &     car\_eval\_4 &    1728 &        21 &                   0 &                    21 &          25.6:1 \\
        20 &   wine\_quality &    4898 &        11 &                  11 &                     0 &          25.8:1 \\
        21 &     letter\_img &   20000 &        16 &                   0 &                    16 &          26.2:1 \\
        22 &      yeast\_me2 &    1484 &         8 &                   6 &                     2 &          28.1:1 \\
        23 &        webpage &   34780 &       300 &                   0 &                   300 &          34.5:1 \\
        24 &    ozone\_level &    2536 &        72 &                  72 &                     0 &          33.7:1 \\
        25 &    mammography &   11183 &         6 &                   6 &                     0 &          42.0:1 \\
        26 &   protein\_homo &  145751 &        74 &                  74 &                     0 &         111.5:1 \\
        27 &     abalone\_19 &    4177 &        10 &                   7 &                     3 &         129.5:1 \\
        \bottomrule
        \end{tabular}
    \caption{Overview of imbalanced datasets involved in the study}
    \label{tab:df_info}
\end{table}

The datasets are shown in table~\ref{tab:df_info}, providing information about the size, number of numerical and categorical features and the imbalance ratio, which is calculated as

\begin{equation}
    \text{imbalance ratio} = \frac{N_{maj}}{N_{min}}
\end{equation}

The datasets are sorted from smallest to biggest ratio between $N_{maj}$ and $N_{min}$. It can be seen that the most balanced dataset has 8.6 samples in the majority class for every minority sample and ranges all the way to 129.5 samples to one. Other variations between datasets include the number of total samples, $N$, ranging from 336 samples in the smallest dataset 'ecoli' to 145.751 samples in the 'protein\_homo' dataset. Larger datasets benefit from the fact that they include more information about each class, whereas small datasets lack the information to discover regularities or patterns in the training data~\cite{Ali2013ClassificationReview}.

A final note is the difference between numerical and categorical features. This thesis labels continuous features as numerical and discrete, ordinal or binary features as categorical. Techniques such as SMOTE, which work on a k-nearest neighbours principle, are meant to be used with numerical features and expect samples in a continuous sample space, while other techniques such as Random Oversampling copy existing samples and therefore do not require specific feature types. These variations in the datasets help in analysis of oversamplers in different scenarios and can help understand where they perform best.

\section{oversamplers}
This thesis tested the performance of ten oversamplers over 27 datasets. Note that SMOTE-NC does not work for exclusively numerical or categorical data, and is only used on the datasets which contain both feature types. The reasoning for SMOTE-NC is that it is an adaption on regular SMOTE, which already works with numerical features. To include a baseline for comparison, the results without oversampling are included as well. Next, five oversamplers from the IMB-learn package were selected. Oversamplers in this package have been published for at least three years with over 200 citations, and have a proven usefulness. These are considered well-established oversamplers and are a good benchmark for Synthsonic. One additional oversampler is added from Kovacs' study~\cite{Kovacs2019AnDatasets}, called Polynomial Fit SMOTE. This oversampler ranked as number one in the study, with the highest results in most metrics. This completes the group of 10 oversamplers to be tested on the imbalanced datasets. 

\section{metrics}
There are already a large variety of widely accepted and standard measures for binary classification, many of these are not suitable when dealing with imbalanced data. Scores such as accuracy are based on the number of correctly classified samples, where the performance of the majority class can be over represented~\cite{Fernandez2018LearningSets}. A classifier which assigns the majority class to all instances will achieve a $99\%$ accuracy if the imbalance ratio is $99:1$. Another drawback is that it is assumed that every error is equally costly. In imbalanced classification it is often more important to correctly label the minority class and misclassifying these cases should be associated with a higher cost. A medical model which fails to predict if a patient has a disease is costlier than a false alarm. Therefore, the goal is to label more minority classes correctly while still maintaining a good score on the majority class. The performance of the classifier is measured as either predicting the correct class labels or associated probability that a sample belongs to a class. Correctly predicting the class labels can be visualized by a confusion matrix, representing the number of TP (true positive), TN (True negative), FP (false positive) and FN (false negative) samples, where $P = TP + FN$ and $N = TN + FP$. The following metrics are then derived from these five measures.

\textbf{G-mean:} the geometric mean of accuracy achieved on minority and majority instances
\begin{equation}
    G = \sqrt{\frac{TP}{P} \cdot \frac{TN}{N}}
\end{equation}

\textbf{Precision:} the ratio between the true positives and false positives. Intuitively it is the ability to not mislabel a sample in the positive class
\begin{equation}
    PR = \frac{TP}{TP + FP}
\end{equation}

\textbf{Recall:} the ratio between true positives and false negatives, or the ability to find all positive samples.
\begin{equation}
    RE = \frac{TP}{TP + FN}
\end{equation}

\textbf{$F_1$-score:} the harmonic mean of precision (PR) and recall (RE)
\begin{equation}
    F_1 = 2 \cdot \frac{PR \cdot RE}{PR + RE}
\end{equation}

\textbf{Balanced accuracy:} the average of recall obtained on each class, it is the accuracy score with class-balanced sample weights
\begin{equation}
    \frac{\frac{TP}{TP + FN} + \frac{TN}{TN + FP}}{2}
\end{equation}

\textbf{PR AUC:} the Area Under Precision-Recall Curve charactarizes the trade-off between precision and recall for different probability thresholds.

\section{Evaluation methods}
cross validation: Due to the small dataset sizes $N$ and lower amount of samples in $N_{min}$, using a single split between training and testing data can lead to a biased result depending on the split. To prevent overfitting, the oversamplers are evaluated over a stratified k-fold cross-validation using four folds to achieve reproducible and reliable results. This number is chosen based on the smallest dataset to ensure that there are enough samples of the positive class in the testing set, the folds are the same for each oversampler to ensure a fair comparison. First, the minority class in the training set is oversampled, with different sampling proportions, before fitting the classifier. After fitting and predicting, the balanced accuracy, geometric mean, F1 score and AUC of the PR curve are determined on an out of hold validation set. This process is repeated for each fold in the cross-validation and the average is taken for a final score.  %Details of your approach (10-15 pages)
%\input{Chapters/ExperimentsandResults} %Evaluation/testing of your hypothesis (10-15 pages)
%\lhead{Discussion}

Discuss %Evaluation/testing of your hypothesis (5-8 pages)
%\input{Chapters/ConclusionsandFutureWork} % (5-6 pages)
\clearpage

\chapter{Results}
\lhead{\emph{Results}}

Ten oversamplers, including Synthsonic, were used in classifying 27 imbalanced datasets shown earlier in table~\ref{tab:df_info}. Due to a large amount of results, it is not feasible to discuss all details of oversampling. A selection of results was made which gives insight into oversampling in general, and the performance of Synthsonic in particular. For comparison, the average results are shown for all datasets, but also for different subsets of datasets based on their characteristics, these will be introduced later in this section.

\section{The proportion parameter}
Each metric is the highest achieved by the oversampler on a dataset. Because all parameters were kept the same except for the proportion, it's therefore checked if certain oversamplers favour a certain proportion. To check this, the amounts were counted which achieved the highest score on all datasets, and the results are shown in table~\ref{tab:prop}. Overall, there is no preferred oversampler across all oversamplers, as all proportions could produce a highest score. Interesting to note is that a 1:1 ratio is not always the best for performance. In earlier sections it was described that classifiers prefer balanced classes, but it seems that evenly balanced classes are not always necessary. The results show that the best performance can also be achieved with a small oversampling proportion. Additionally, ADASYN and BorderlineSMOTE seem to perform better at lower proportions, both achieving the most high-scores with a 0.2 proportion. The other oversamplers achieve highest use a variety of all proportions to achieve the best scores. Most notable is SVMSMOTE, which achieved highest scores with all proportions except 0.4. 

The effect of the oversampling proportion was further studied by looking at the average proportion that produced the highest scores per dataset. 

\begin{table}[]
\centering
\resizebox{\textwidth}{!}{
\begin{tabular}{l|rrrrrrrrr}
    \toprule
    proportion &  ADASYN &  BorderlineSMOTE &  RandomOversampler &  Random\_SMOTE &  SMOTE &  SMOTENC &  SVMSMOTE &  polynom\_fit\_SMOTE &  synthsonic \\
    \midrule
    0.2 &       9 &                7 &                  6 &             3 &      7 &        1 &         6 &                  5 &           7 \\
    0.4  &       4 &                5 &                  5 &             6 &      2 &        4 &         3 &                  4 &           6 \\
    0.6 &       7 &                3 &                  4 &             4 &      6 &        1 &         6 &                  7 &           5 \\
    0.8 &       4 &                6 &                  4 &             6 &      5 &        3 &         6 &                  4 &           2 \\
    1.0 &       3 &                6 &                  8 &             8 &      7 &        2 &         6 &                  7 &           7 \\
    \bottomrule
\end{tabular}}
\caption{Amount of highest f1-scores achieved per oversampler by proportion.}
\label{tab:prop}
\end{table}


\section{Ranking of oversamplers}
In order to create a fair ranking over all datasets, the highest result per metrics was taken when using XGBoost. Table~\ref{tab:ranking} provides an overview of the performance on all datasets. ADASYN performs best on balanced accuracy and geometric mean, while SVMSMOTE has the highest f1 score and Synthsonic has the highest PR-AUC. A surprise is that Polynom\_fit\_SMOTE only reaches the top five in PR-AUC, as it was the best performing oversampler in the study of Kovacs et al.~\cite{Kovacs2019AnDatasets}. Additionally, ADASYN did not reach the top 10 in any metric in their study, while here it outperforms the other oversamplers. Lowest scores in table~\ref{tab:ranking} are achieved by the baseline performance and SMOTENC. SMOTENC performs significantly worse than other oversamplers, and even below the baseline performance on f1 score and PR-AUC. Overall, the differences between oversamplers is little. It shows that the original SMOTE algorithm performs well and that the adaptations do not guarantee a higher performance.

\begin{table}[ht]
\centering
    \resizebox{\textwidth}{!}{
\begin{tabular}{llrlrlrlr}
\toprule
{} &        Oversampler &  balanced\_accuracy &        Oversampler &  G\_mean &        Oversampler &      f1 &        Oversampler &  pr\_auc \\
\midrule
1  &             ADASYN &             \textbf{0.7983} & ADASYN &  \textbf{0.7296} & SVMSMOTE &  \textbf{0.6262} & synthsonic &  \textbf{0.6494} \\
2  &              SMOTE &             0.7979 &              SMOTE &  0.7294 &             ADASYN &  0.6238 &           SVMSMOTE &  0.6494 \\
3  &         synthsonic &             0.7973 &         synthsonic &  0.7268 &              SMOTE &  0.6220 &             ADASYN &  0.6492 \\
4  &           SVMSMOTE &             0.7966 &  RandomOversampler &  0.7263 &       Random\_SMOTE &  0.6217 &  RandomOversampler &  0.6486 \\
5  &  RandomOversampler &             0.7963 &       Random\_SMOTE &  0.7256 &    BorderlineSMOTE &  0.6205 &  polynom\_fit\_SMOTE &  0.6471 \\
6  &       Random\_SMOTE &             0.7959 &           SVMSMOTE &  0.7249 &  RandomOversampler &  0.6201 &              SMOTE &  0.6471 \\
7  &    BorderlineSMOTE &             0.7942 &    BorderlineSMOTE &  0.7242 &         synthsonic &  0.6141 &       Random\_SMOTE &  0.6456 \\
8  &  polynom\_fit\_SMOTE &             0.7820 &  polynom\_fit\_SMOTE &  0.6980 &  polynom\_fit\_SMOTE &  0.6063 &    BorderlineSMOTE &  0.6429 \\
9  &            SMOTENC &             0.7681 &            SMOTENC &  0.6934 &     NoOversampling &  0.5804 &     NoOversampling &  0.6419 \\
10 &     NoOversampling &             0.7617 &     NoOversampling &  0.6536 &            SMOTENC &  0.5553 &            SMOTENC &  0.5711 \\
\bottomrule
\end{tabular}}
\caption{A ranking of oversamplers based on their average score of all datasets using XGBoost.}
\label{tab:ranking}
\end{table}

In addition, the reader can see from figure~\ref{fig:Total} that all oversamplers perform similarly, and the differences are minor. 

\begin{figure}[htp]
\centering
    \includesvg[width=.45\textwidth]{../Plots/Total/total_balanced_accuracy.svg}\quad
    \includesvg[width=.45\textwidth]{../Plots/Total/total_G_mean.svg}
    \medskip
    \includesvg[width=.45\textwidth]{../Plots/Total/total_f1.svg} \quad
    \includesvg[width=.45\textwidth]{../Plots/Total/total_pr_auc.svg}
    \medskip
    \includesvg[width=.45\textwidth]{../Plots/Total/total_precision.svg}\quad
    \includesvg[width=.45\textwidth]{../Plots/Total/total_recall.svg}
\caption{Average metrics scores per oversampler on all datasets. All scores range from 0 (worst) to 1 (best).}
\label{fig:Total}
\end{figure}

\section{Performance with K-nearest Neighbours classifier}
The same experiments are also performed using a K-nearest Neighbours classifier with 3 and 5 nearest neighbours to validate the results. According to Kovacs et al.~\cite{Kovacs2019AnDatasets}, oversampling has a larger benefit with simpler models. Using XGBoost as a baseline performance resulted in minimal returns from oversampling, but this is not the case for KNN. Table ~\ref{tab:ranking_knn} shows that the difference between the baseline performance and oversampling is significant for balanced accuracy and geometric mean, but the difference in f1 and PR-AUC is smaller. These results do confirm that the performance between oversamplers is still similar, but the rankings have changed. ADASYN no longer ranks the highest in balanced accuracy and Geometric mean and drops to second place, while the Geometric mean score improves by 5,6\%. The baseline performance now shows that oversampling significantly improves the balanced accuracy and Geometric mean, with a similar relative improvement in f1 score as with XGBoost. Note that XGBoost scored a higher f1 score than k-nearest neighbours. The largest difference between the two classifiers is in the PR-AUC. Not only are the scores lower in table~\ref{tab:ranking_knn}, but oversampling performs worse than the baseline performance.

\begin{table}[ht]
\centering
    \resizebox{\textwidth}{!}{
\begin{tabular}{llrlrlrlr}
\toprule
{} &        Oversampler &  balanced\_accuracy &        Oversampler &  G\_mean &        Oversampler &      f1 &        Oversampler &  pr\_auc \\
\midrule
1 &       Random\_SMOTE & \textbf{0.7917} & Random\_SMOTE &  \textbf{0.7705} &  RandomOverSampler &  \textbf{0.4692} & NoOversampler & \textbf{0.5166} \\
2 &             ADASYN &             0.7882 &             ADASYN &  0.7656 &         synthsonic &  0.4663 &       Random\_SMOTE &  0.5099 \\
3 &              SMOTE &             0.7861 &              SMOTE &  0.7628 &           SVMSMOTE &  0.4659 &  polynom\_fit\_SMOTE &  0.5084 \\
4 &  polynom\_fit\_SMOTE &             0.7793 &    BorderlineSMOTE &  0.7384 &    BorderlineSMOTE &  0.4632 &    BorderlineSMOTE &  0.5054 \\
5 &    BorderlineSMOTE &             0.7768 &  polynom\_fit\_SMOTE &  0.7351 &              SMOTE &  0.4481 &           SVMSMOTE &  0.5038 \\
6 &           SVMSMOTE &             0.7677 &           SVMSMOTE &  0.7212 &       Random\_SMOTE &  0.4433 &              SMOTE &  0.5031 \\
7 &         synthsonic &             0.7586 &         synthsonic &  0.7024 &             ADASYN &  0.4419 &         synthsonic &  0.5010 \\
8 &  RandomOverSampler &             0.7512 &  RandomOverSampler &  0.6956 &  polynom\_fit\_SMOTE &  0.4313 &             ADASYN &  0.5005 \\
9 &      NoOversampler &             0.6693 &      NoOversampler &  0.4794 &      NoOversampler &  0.4043 &  RandomOverSampler &  0.4885 \\
\bottomrule
\end{tabular}}
\caption{A ranking of oversamplers based on their average score of all datasets using K-nearest neigbhours (k=5).}
\label{tab:ranking_knn}
\end{table}

The difference in f1 score and PR-AUC is caused by the shift in precision and recall, as shown in table~\ref{tab:prrecknn}. The results using XGBoost showed that recall increases due to oversampling, with a slight decrease in precision. This effect is amplified when K-nearest neighbours is used. The precision score drops by 40\% from the baseline performance to the lowest performer, polynom\_fit\_SMOTE, while the recall score is nearly doubled. The discrepancy between precision and recall affects the f1 score, the harmonic mean of precision and recall, and the area under the PR-curve. As mentioned earlier, the goal is to increase sensitivity to the minority class, without losing performance on the majority class.

\begin{table}[h]
\centering
\resizebox{0.5\textwidth}{!}{
\begin{tabular}{lrrrr}
\toprule
{} &  precision &  precision\_std &  recall &  recall\_std \\
Oversampler       &            &                &         &             \\
\midrule
NoOversampler     &     0.5983 &         0.1108 &  0.3460 &      0.0445 \\
synthsonic        &     0.4329 &         0.0320 &  0.5819 &      0.0615 \\
RandomOverSampler &     0.4267 &         0.0375 &  0.5672 &      0.0582 \\
SVMSMOTE          &     0.4031 &         0.0326 &  0.6120 &      0.0589 \\
BorderlineSMOTE   &     0.3899 &         0.0305 &  0.6450 &      0.0570 \\
SMOTE             &     0.3688 &         0.0248 &  0.6843 &      0.0596 \\
ADASYN            &     0.3602 &         0.0265 &  0.6943 &      0.0617 \\
Random\_SMOTE      &     0.3584 &         0.0243 &  0.7105 &      0.0598 \\
polynom\_fit\_SMOTE &     0.3556 &         0.0261 &  0.6876 &      0.0466 \\
\bottomrule
\end{tabular}}
    \caption{Average precision and recall scores for K-nearest neighbours classifier on all datasets ($k=5$).}
    \label{tab:prrecknn}
\end{table}

\newpage
The experiment was repeated with a K-nearest neighbours classifier with $k=3$, results in table~\ref{tab:rankingknn3} are similar to the results of $k=5$. The baseline performance is improved for all metrics, while the oversampling performance does not significantly change. The f1 score improves but is still below the performance of XGBoost. While oversampling does seem like a valid option for classifier improvement, the same effect can be achieved by selecting a better classifyer.

\begin{table}[h]
\centering
\resizebox{\textwidth}{!}{
\begin{tabular}{llrlrlrlr}
\toprule
{} &        Oversampler &  balanced\_accuracy &        Oversampler &  G\_mean &        Oversampler &      f1 &        Oversampler &  pr\_auc \\
\midrule
1 &       Random\_SMOTE & \textbf{0.7852} &       Random\_SMOTE &  \textbf{0.7613} & SVMSMOTE &  \textbf{0.4712} & Random\_SMOTE &  \textbf{0.5112} \\
2 &             ADASYN &             0.7823 &             ADASYN &  0.7566 &  RandomOverSampler &  0.4687 &  polynom\_fit\_SMOTE &  0.5094 \\
3 &              SMOTE &             0.7792 &              SMOTE &  0.7521 &    BorderlineSMOTE &  0.4683 &    BorderlineSMOTE &  0.5066 \\
4 &  polynom\_fit\_SMOTE &             0.7770 &  polynom\_fit\_SMOTE &  0.7347 &         synthsonic &  0.4625 &              SMOTE &  0.5032 \\
5 &    BorderlineSMOTE &             0.7693 &    BorderlineSMOTE &  0.7262 &              SMOTE &  0.4511 &           SVMSMOTE &  0.5024 \\
6 &           SVMSMOTE &             0.7622 &           SVMSMOTE &  0.7123 &       Random\_SMOTE &  0.4468 &             ADASYN &  0.5023 \\
7 &         synthsonic &             0.7535 &         synthsonic &  0.6913 &             ADASYN &  0.4463 &      NoOversampler &  0.5013 \\
8 &  RandomOverSampler &             0.7351 &  RandomOverSampler &  0.6664 &  polynom\_fit\_SMOTE &  0.4403 &         synthsonic &  0.4991 \\
9 &      NoOversampler &             0.6809 &      NoOversampler &  0.5216 &      NoOversampler &  0.4292 &  RandomOverSampler &  0.4838 \\
\bottomrule
\end{tabular}}
    \caption{Caption}
    \label{tab:rankingknn3}
\end{table}


\section{Top oversamplers on different dataset types}
The datasets in this study are divided according to their characteristics: by imbalance ratio (IR), number of minority class samples, number of features and feature type (numerical or categorical). These characteristics were highlighted in table~\ref{tab:df_info}. Thresholds are introduced to separate low ($IR < 20$) and high ($IR > 20$) IR, low ($N_{min} < 100$) and high ($N_{min} > 100$) number of minority samples, low  ($f < 40$)  and high ($f > 40$) number of features, numerical and categorical features, and low ($N < 1.000$) and high ($N > 1.000$) sample size. These thresholds are empirical and chosen for these specific datasets. For example, Alcala et al.~\cite{Alcala-Fdez2011KEELFramework} classify a dataset as highly imbalanced when $IR > 9$. This threshold would not provide any insights in this study, as it would classify all but one dataset as highly imbalanced. Therefore, thresholds are chosen so that they give insight in the performance on these particular datasets.

From the overall results in table~\ref{tab:ranking}, it became clear that the difference in scores between oversamplers is minor. Therefore, results are handled differently here. Instead of reporting the scores, oversamplers are ranked from $1^{st}$ to $10^{th}$ place based on the average of their performance. First, the average score per metric is calculated and are then ranked from best to worst. This is repeated for each metric, so that an oversampler has a rank for all metrics. These ranks are then averaged to produce a final average rank per subset of datasets. Rankings are easier to interpret than raw scores and allow the reader to quickly see what an oversampler's best performance is. The final results are shown in table~\ref{tab:subset_performance} 

The first thing to observe is that there is no oversampler which manages to rank the highest on all subsets. Next, the baseline performance is often ranked among the lowest, but still outperformed some oversamplers on seven subsets. On two of those, it shared the last place with SMOTE-NC. Surprisingly, SMOTE-NC is often ranked the lowest but still achieved the a third and first place on low IR and high features datasets, making it a valid option on unseen datasets with these characteristics. Alternatively, ADASYN ranks highest on high IR, low $N_{min}$, numerical and low size datasets. Another thing to note is that SMOTE performs well on all subsets. It does not rank first, but consistently outperforms other SMOTE-variants. Even on categorical datasets, where SMOTE is not built for, it ranks second. Together with the results from table~\ref{tab:ranking} and figure~\ref{fig:Total}, it shows that regular SMOTE is a valid technique on various dataset types, and the benefit of SMOTE-variants is minor. Having said that, ADASYN performs relatively better, both in the total results as in the different subsets.

Table~\ref{tab:subset_performance} also shows where Synthsonic outperforms other oversamplers. First, Synthsonic performs significantly better on datasets with a low IR, meaning that the classes are more balanced. Ranking first in low IR datasets does not translate to the same performance on more imbalanced datasets, as Synthsonic drops down to the sixth rank. For low $X_{min}$, low features and low size, synthsonic performs better than the higher alternatives. In other words, Synthsonic ranks higher on smaller datasets with fewer features. The Bayesian Network in Synthsonic attributes to the difference in performance on low and high feature datasets, as it captures the dependencies between all features. If the model is not optimal, or when dependencies are hard to find, the performance of Synthsonic will suffer. As the number of features increases, it also becomes harder to find an accurate network of the features. 

\begin{table}[]
\centering
\resizebox{\textwidth}{!}{
    \begin{tabular}{lrrrrrrrrrr}
    \toprule
    {} &  low IR &  high IR &  low $X_{min}$ &  high $X_{min}$ &  low features &  high features &  num &  cat &  low size &  high size \\
    oversampler & $IR < 20$ & $IR > 20$ & $N_{min} < 100$  & $N_{min} > 100$ & $f < 40$ & $f > 40$ &          &      &  $N < 1.000$ & $N > 1.000$\\
    \midrule
    ADASYN            &       6 &        1 &            1 &             5 &             4 &              3 &    1 &    3 &         1 &          3 \\
    BorderlineSMOTE   &       9 &        3 &            6 &             4 &             7 &              6 &    4 &    8 &         6 &          7 \\
    No Oversampling    &      10 &        9 &           10 &             8 &             9 &             10 &    9 &    9 &         9 &          8 \\
    RandomOversampler &       4 &        6 &            7 &             2 &             2 &              7 &    7 &    1 &         8 &          1 \\
    Random\_SMOTE      &       7 &        4 &            4 &             7 &             6 &              5 &    6 &    3 &         3 &          6 \\
    SMOTE             &       5 &        2 &            2 &             3 &             5 &              2 &    3 &    2 &         3 &          4 \\
    SMOTENC           &       3 &       10 &            9 &            10 &            10 &              1 &    9 &   10 &         9 &         10 \\
    SVMSMOTE          &       2 &        5 &            5 &             1 &             3 &              4 &    2 &    6 &         6 &          2 \\
    Polynom\_fit\_SMOTE &       8 &        8 &            8 &             9 &             7 &              9 &    8 &    7 &         5 &          9 \\
    Synthsonic        &       1 &        6 &            3 &             5 &             1 &              8 &    5 &    5 &         2 &          5 \\
    \bottomrule
    \end{tabular}}
\caption{Ranking of oversamplers on different dataset types. The rank is averaged over all metrics for the datasets in that subset.}
\label{tab:subset_performance}
\end{table}

\section{Runtime}
The average runtime over all 27 datasets is shown in table~\ref{tab:runtimes}. These results are achieved by oversampling each dataset using a sampling strategy of 1, which means both classes contain an equal amount of samples. From the average runtimes it can be seen that the SMOTE-based oversamplers all have a comparable average runtime. The SMOTE-based oversamplers all average under 1 second for oversampling, with the exception of SMOTENC and SVMSMOTE. These oversamplers reached an average runtime of 1.395 seconds and 2.725 seconds respectively. From table~\ref{tab:runtimes} it can be seen that Synthsonic differs significantly from its counterparts. The average runtime was 86 seconds. This is largely due to a few outliers which result in exceptionally high runtimes. These peaks are also present in other oversamplers, as seen in figure~\ref{fig:runvfeat}. SMOTE-based oversamplers also suffer a hit in runtime, but are still able to stay far under the runtime of Synthsonic. Lastly, figure~\ref{fig:runvfeat} shows that the runtime of Synthsonic is affected by the amount of features in the dataset.

Next, take a look at the spikes in runtime for SMOTE-based oversamplers in figure~\ref{fig:runvfeat}. While Synthsonic's runtime increases with the number of features, other oversamplers maintain a similar runtime over the datasets, with the exception of two big spikes. Also note that the runtimes are on a different scale, but all follow a similar pattern. The dataset responsible for the first big spike in runtime is called 'protein\_homo', the largest dataset in this study. Figure~\ref{fig:runvsize} shows the relation between the dataset size and the oversamplers' runtime, where the runtime rapidly increases for large datasets. However, keep in mind that the increase in runtime is significant on a relative scale, as the oversampling runtime is still under 5 seconds. 

\begin{table}[h]
    \centering
    \begin{tabular}{lrr}
    \toprule
    oversampler & mean run-  &  stdev \\
     & time (s) & \\
    \midrule
    RandomOversampler &    0.014 &        0.001 \\
    polynom\_fit\_SMOTE &    0.041 &        0.002 \\
    SMOTE             &    0.059 &        0.002 \\
    BorderlineSMOTE   &    0.297 &        0.005 \\
    ADASYN            &    0.397 &        0.006 \\
    Random\_SMOTE      &    0.419 &        0.015 \\
    SMOTENC           &    1.395 &        0.024 \\
    SVMSMOTE          &    2.735 &        0.057 \\
    synthsonic        &   86.641 &        0.646 \\
    \bottomrule
    \end{tabular}
    \caption{Runtime for 1:1 oversampling, averaged over all 27 datasets.}
    \label{tab:runtimes}
\end{table}

\begin{figure}[h!]
\centering
\begin{subfigure}[b]{0.7\textwidth}
    \centering
    \includesvg[width=\textwidth]{../Plots/Runtime/boxplot.svg}
    \caption{A boxplot of the runtimes per oversampler over all datasets, oversampled to a ratio of 1:1.}
    \label{fig:boxplot}
\end{subfigure}
\vfill
\begin{subfigure}[b]{0.7\textwidth}
    \centering
    \includesvg[width=\textwidth]{../Plots/Runtime/runtime_vs_features.svg}
    \caption{The runtime for 1:1 oversampling, sorted from smallest to largest dataset.}
    \label{fig:runvfeat}
\end{subfigure}
\vfill
\begin{subfigure}[b]{0.7\textwidth}
    \centering
    \includesvg[width=\textwidth]{../Plots/Runtime/runtime_vs_size.svg}
    \caption{he runtime for 1:1 oversampling, sorted from fewest to most features.}
    \label{fig:runvsize}
\end{subfigure}

\end{figure}




%----------------------------------------------------------------------------------------
%	THESIS CONTENT - APPENDICES
%----------------------------------------------------------------------------------------

%\addtocontents{toc}{\vspace{2em}} % Add a gap in the Contents, for aesthetics

%\appendix % Cue to tell LaTeX that the following 'chapters' are Appendices

% Include the appendices of the thesis as separate files from the Appendices folder
% Uncomment the lines as you write the Appendices

%\input{Appendices/AppendixA}
%\input{Appendices/AppendixB}
%\input{Appendices/AppendixC}

\addtocontents{toc}{\vspace{2em}} % Add a gap in the Contents, for aesthetics

\backmatter

%----------------------------------------------------------------------------------------
%	BIBLIOGRAPHY
%----------------------------------------------------------------------------------------

\label{References}

\lhead{\emph{References}} % Change the page header to say "Bibliography"

\bibliographystyle{unsrtnat} % Use the "unsrtnat" BibTeX style for formatting the Bibliography

\bibliography{References.bib} % The references (bibliography) information are stored in the file named "Bibliography.bib"

\end{document}  